% Paquetes esenciales
\usepackage[utf8]{inputenc}
\usepackage[T1]{fontenc}
\usepackage[english]{babel}
\usepackage{amsmath, amssymb, amsfonts, amsthm}
\newtheorem{proposition}{Proposition}
\newtheorem{definition}{Definition}
\usepackage{graphicx}
\usepackage{hyperref}
\usepackage{xcolor}
\usepackage{geometry}
\usepackage{fancyhdr}
\usepackage{setspace}
\usepackage[most]{tcolorbox}
\usepackage{listings}
\usepackage{tcolorbox}
\tcbuselibrary{breakable, skins}
\usepackage{minted} % Asegúrate de compilar con -shell-escape
\usepackage{etoolbox} 
\usepackage{xparse} % permite definir comandos con 
\usepackage[round,authoryear]{natbib}
\usepackage[ruled,vlined,linesnumbered]{algorithm2e}
\usepackage{booktabs}

% Code blocks
\definecolor{codegray}{gray}{0.95}
\lstset{
  backgroundcolor=\color{codegray},
  basicstyle=\ttfamily\small,
  frame=single,
  numbers=left,
  numberstyle=\tiny,
  tabsize=2,
  breaklines=true,
  showstringspaces=false
}

% Hipervínculos
\hypersetup{
    colorlinks=true,
    linkcolor=blue,
    urlcolor=blue,
    citecolor=blue
}

% \newtcolorbox{answerbox}{
%   colback=white,
%   colframe=black,
%   boxrule=0.8pt,
%   arc=3pt,
%   left=5pt,
%   right=5pt,
%   top=5pt,
%   bottom=5pt
% }

% \newcommand{\answer}[1]{\begin{answerbox}#1\end{answerbox}}
